% ---- pg-445 （书页 433） ----

例如, 第一样本典型变量配对为
\[
\begin{aligned}
\hat{U}_1&=0.42z_1^{(1)}+0.21z_2^{(1)}+0.17z_3^{(1)}-0.02z_4^{(1)}+0.44z_5^{(1)}\\
\hat{V}_1&=0.42z_1^{(2)}+0.22z_2^{(2)}-0.03z_3^{(2)}+0.01z_4^{(2)}+0.29z_5^{(2)}+0.52z_6^{(2)}-0.12z_7^{(2)}
\end{aligned}
\]
且样本典型相关系数为 $\hat{\rho}_1^*=0.55$.

根据系数, $\hat{U}_1$ 首先是一个用户反馈和自主权变量, 而 $\hat{V}_1$ 则表示对主管、事业前景和工种的满意度, 以及公司认同感.

为提供对 $\hat{U}_1$ 和 $\hat{V}_1$ 的解释, 计算 $\hat{U}_1$ 与其成分变量之间的样本相关系数和 $\hat{V}_1$ 与其成分变量之间的样本相关系数. 下表列出了一个组中变量与另一个组的第一样本典型变量的样本相关系数. 这些相关系数可由\hyperref[eq:10-29]{式(10-29)}计算得到.

\begin{mvtable}
\mvtabcap{原始变量与典型变量的样本相关系数}
\begin{tabular}{l|cc|l|cc}
\hline
\multirow{2}{*}{$\mathbf{X}^{(1)}$ 变量} & \multicolumn{2}{c|}{样本典型变量} & \multirow{2}{*}{$\mathbf{X}^{(2)}$ 变量} & \multicolumn{2}{c}{样本典型变量}\\
\cline{2-3}\cline{5-6}
& $\hat{U}_1$ & $\hat{V}_1$ & & $\hat{U}_1$ & $\hat{V}_1$\\
\hline
1. 用户反馈 & 0.83 & 0.46 & 1. 主管满意度 & 0.42 & 0.75\\
2. 任务重要性 & 0.74 & 0.41 & 2. 事业前景满意度 & 0.35 & 0.65\\
3. 任务多样性 & 0.75 & 0.42 & 3. 财政满意度 & 0.21 & 0.39\\
4. 任务特性 & 0.62 & 0.34 & 4. 工作强度满意度 & 0.21 & 0.37\\
5. 自主权 & 0.85 & 0.48 & 5. 公司认同感 & 0.36 & 0.65\\
& & & 6. 工种满意度 & 0.44 & 0.80\\
& & & 7. 总体满意度 & 0.28 & 0.50\\
\hline
\end{tabular}
\end{mvtable}

所有五个表示职业特性的变量具有与第一典型变量 $\hat{U}_1$ 大致相同的相关系数. 从此点看, $\hat{U}_1$ 可以形容为职业特性“指标”. 它与我们喜欢的基于系数的解释不同, 这里任务变量并不重要.

第一典型变量对中的另一个成员 $\hat{V}_1$, 看起来主要地代表了主管满意度、事业前景满意度、公司认同感和工种满意度. 正像变量所表达的, $\hat{V}_1$ 可以视为职业满意度-公司认同感指标. 这些与前述的基于诸 $z_i^{(2)}$ 典型系数的解释相一致. 两个指标 $\hat{U}_1$ 和 $\hat{V}_1$ 间的样本相关系数为 $\hat{\rho}_1^*=0.55$. 看来, 职业特性和职业满意度之间有部分重叠, 我们将在例 10.7 中进一步考察这个问题.\hfill$\blacksquare$

第一对 $(\hat{U}_1,\hat{V}_1)$ 的散布图可以揭示非典型的、需要进一步研究的观察值 $\mathbf{x}_j$. 若典型相关系数 $\hat{\rho}_2^*,\hat{\rho}_3^*,\cdots$ 亦较大, $(\hat{U}_2,\hat{V}_2),(\hat{U}_3,\hat{V}_3),\cdots$ 的散布图在这方面还是有用的. 作为阐述实际题材的一个辅助, 很多分析者建议标绘“重要的”典型变量而除去它们的成分变量. 这些图加强了\hyperref[eq:10-29]{式(10-29)}中的相关系数.

若样本较大, 我们常常希望将其分为两半, 前一半可用来构造并评价该样本的典型变量和典型相关系数, 其结果用剩下的观察值来“验证”. 典型分析中的实质性改变(如果有的话)会为取样的不定性和结论的稳定性提供一个提示.
